\subsection{Recursive fibers, target paths and epistemic cuts}
An unresolved transformation is
\begin{equation}f=(o_f,q_f,\gamma_f,r_f,\operatorname{parent}(f)).\end{equation}
Residual routing opens only dimensions capable of changing the target conclusion. For target $\tau=(q,\alpha^*,\gamma)$, \RAKL{} searches for an authority-valid support subgraph or hypergraph. A conceptual low-cost cut objective is
\begin{equation}B^*_{\tau}=\arg\min_B\operatorname{Cost}(B)\end{equation}
subject to every admissible support route intersecting $B$. This equation defines a research-control target rather than a guarantee that a globally optimal cut always exists or is computationally tractable. A cut may correspond to missing evidence, measurement, context, transition, mechanism, formal lemma or reasoning operator.

\subsection{Epistemic cuts turn ``what is missing?'' into a constructive object}

A research system often knows that it cannot support a target but not why. The epistemic-cut formulation makes the missing support explicit. Consider the hypergraph of admissible support routes from current evidence and assumptions to a target authority certificate. A cut is a set of missing or blocking obligations whose resolution would intersect every currently viable route. Depending on the problem, a cut can contain a measurement, a calibration, a context coordinate, a causal assumption, a formal lemma, a replication, a transformation witness or a method operator.

This is not claimed as a new graph-theoretic min-cut theorem. The scientific content lies in the typed elements and their cost. A cheap new measurement and an expensive randomized intervention may both separate the same mechanisms but carry different feasibility, ethical and informational profiles. The cut representation provides a common interface for experiment design and method diagnosis while preserving those types.

Cuts also prevent premature invention. If the target is blocked because a raw measurement is absent, inventing a new theory does not close the gap. If every route depends on an unidentified nuisance transformation, collecting more of the same outcome may not help. If the missing object is a method operator, ordinary evidence acquisition cannot solve it. The cut therefore routes the next action to the correct layer.

When multiple cuts exist, the system can optimize over cost, expected target change and risk, but the optimization criterion must be registered. A lowest monetary-cost cut may be scientifically weak; a high-information experiment may be infeasible; a method invention may be cheap to propose but expensive to validate. The action-selection layer is deliberately downstream of cut typing so it cannot compare incomparable actions without an explicit utility model.

The concept also clarifies saturation. A fiber with a material unresolved cut is not saturated merely because search rounds are flat. Flat search plus an explicit unavailable measurement can be a terminal block; flat search plus an unattempted feasible discriminator is not. Saturation is therefore a statement about the status of remaining cuts as much as about novelty counts.

\subsection{Action selection, model criticism and assumption sensitivity}
When calibrated probabilities are justified, one admissible policy family is
\begin{equation}u(a\mid K_t)=\frac{\lambda_Q I(Q;Y_a\mid K_t)+\lambda_M\operatorname{Sep}(a,\mathcal V_t)+\lambda_N E[\Delta N_a]}{\operatorname{Cost}(a)}.\end{equation}
Without calibrated probabilities, \RAKL{} uses set-valued alternatives such as worst-case mechanism separation or identified-set shrinkage.

For frozen model criticism, scientific discrepancy probes are $s_k=T_k(y_{\mathrm{obs}})$ and are compared with predictive/generative samples $s_k^{(b)}=T_k(y_{\mathrm{pred}}^{(b)})$. A two-sided finite-sample predictive tail plus a predeclared materiality tolerance can identify a structured residual. Passing the probes establishes adequacy only relative to the frozen probe family, population and resolution; it does not prove truth or mechanism identity.

A complementary assumption-sensitivity operator evaluates a frozen envelope $(a_j,\theta_j)$ under the same context and QoI. For material threshold $\delta$,
\begin{equation}
C(\theta;\delta)=
\begin{cases}
\mathrm{POSITIVE}, & \theta>\delta,\\
\mathrm{NEGATIVE}, & \theta<-\delta,\\
\mathrm{INDETERMINATE}, & |\theta|\le\delta.
\end{cases}
\end{equation}
A baseline conclusion is robust only within the registered envelope when all available preregistered scenarios preserve its material conclusion class. Robustness does not prove the assumptions true; a changed population or QoI is a new contextual chart rather than an assumption perturbation of the same claim.

\subsection{Epistemic state as an event-sourced scientific object}

The tuple notation above is compact, but its operational interpretation is closer to an event-sourced state machine than to a mutable notebook. Canonical state is reconstructed from addressable evidence objects and licensed transition events. Each event records the object identities it reads, the evidence and context it relies on, the certificates it creates or withdraws, and the governance identity that authorized the operation. The purpose is not software fashion. Event history is what permits a later reviewer to answer questions that ordinary prose obscures: Which observation first changed the mechanism survivor set? Which assumption made two claims comparable? Which refutation closed a route? Which method version generated a result?

This produces four invariants. \textbf{Identity invariance} requires a frozen object identifier to continue denoting the same object definition; changed scope produces a successor, not retroactive mutation. \textbf{Lineage invariance} requires a derived claim to retain a path to its evidence roots. \textbf{Authority invariance under pure representation operations} requires that normalization, indexing, compression or visualization cannot raise authority. \textbf{Negative-history persistence} requires refutations and failed candidates to remain addressable. Scientific conclusions can still change. What is prohibited is changing the historical record in order to make the current conclusion appear inevitable.

The resulting state transition can be decomposed as
\[
(K_t,p,e,\gamma,g)
\overset{\mathrm{validate}}{\longrightarrow}
v
\overset{\mathrm{authorize}}{\longrightarrow}
u
\overset{\mathrm{commit}}{\longrightarrow}
K_{t+1},
\]
where \(p\) is a proposal, \(e\) evidence, \(\gamma\) context and \(g\) governance state. Validation establishes a typed outcome; authorization checks whether that outcome licenses the proposed authority effect; commit creates the new versioned state. A failure at any stage is retained with a reason. This decomposition prevents an implementation success such as ``the parser produced a result'' from being mistaken for a scientific success such as ``the evidence supports a mechanism.''

A useful consequence is that \CannotCheck{} is a first-class terminal state for an operation. In ordinary chat interfaces, unanswered questions invite guessing; in scientific software, missing calibration metadata, unavailable raw observations, unresolved source identity or insufficient evaluator independence are meaningful states. Returning \CannotCheck{} protects the difference between ``false'' and ``not currently verifiable.'' The same logic motivates \CannotCompile{} when mandatory epistemic context exceeds a model budget: the system must narrow the operation, retrieve a better representation or use a larger resource envelope rather than silently omit a falsifier.

\subsection{Proof obligations travel with transitions}

Every typed transition has positive and negative obligations. Positive obligations state what must be true for the transition to be used: aligned target, compatible units, valid observation map, required evidence, admissible regime and an error law. Negative obligations state what the transition may not create: mechanism authority from association alone, independence from duplicated provenance, identification from one arbitrary representative, or decision authority outside the evaluated loss and population. Writing the negative side is essential because scientific-agent failure often occurs through an implicit coercion rather than an explicitly false equation.

For a transition \(T\), let
\[
\mathcal O(T)=
(\mathcal O_T^+,\mathcal O_T^-,\mathcal F_T,\mathcal R_T)
\]
contain positive obligations, forbidden authority effects, falsifiers and remainder/error semantics. Composition \(T_2\circ T_1\) is licensed only if the output role and context of \(T_1\) meet the input requirements of \(T_2\), the forbidden effects are not crossed, and the combined approximation error is explicitly controlled. A chain of individually reasonable transformations can therefore be blocked globally. This is intentional: most scientific pipelines are reliable only if the semantics of intermediate transformations survive composition.

\subsection{Action selection must distinguish information gain from confirmation pressure}

The action utility
\[
u(a\mid K_t)=
\frac{\lambda_Q I(Q;Y_a\mid K_t)+
\lambda_M\operatorname{Sep}(a,\mathcal V_t)+
\lambda_N E[\Delta N_a]}{\operatorname{Cost}(a)}
\]
is a template, not a universal law. Its terms are meaningful only when their models are meaningful. Expected information gain depends on a probabilistic model for outcomes. Mechanism separation depends on a declared survivor set and observation map. Expected semantic novelty is a planning signal rather than scientific evidence. The weights express governance preferences and should not be tuned after seeing which action makes the favored theory win.

A robust alternative is set based. If calibrated probabilities are unavailable, an action can be valued by worst-case reduction in an identified set, by the minimum pairwise separation it induces among viable mechanisms, or by whether it crosses a blocking proof obligation. This often produces a different choice from selecting the action with the highest expected predictive improvement. The difference is scientifically important when the target is explanation rather than forecast.

Model criticism is a companion to action selection. After a model wins a relative comparison, posterior predictive checks, residual structure, invariance tests, assumption sensitivity and out-of-support diagnostics ask whether the model is adequate for the target. A model can be best in class and still fail adequacy. Conversely, a model can be scientifically adequate for a decision even if a more complex model fits slightly better. The stopping criterion should therefore reference target adequacy and remaining discriminators, not leaderboard rank alone.

Assumption sensitivity turns hidden philosophical disagreement into an empirical or decision object. If a conclusion flips under small plausible changes to an unverified assumption, the authority certificate should expose that sensitivity. If the conclusion is robust across a large admissible assumption set, the system can grant stronger decision authority without pretending that the assumption was observed. This is especially useful for causal and partially identified problems, where the strongest honest result may be a region of conclusions stable across assumption worlds.

Finally, portfolio control prevents local information gain from consuming the entire research budget. Some capacity is reserved for replication, adversarial search, ontology escape, high-risk discriminators and method reflection. This is analogous to the workspace reservation rule: relevance-ranked exploitation is useful but can entrench the current ontology if it monopolizes attention.
